% ============================================================
% 6. Edge Cases \& Structural Extremities (Cardinality)
% ============================================================
\section{Edge Cases \& Structural Extremities: Cardinality Breakdown}
\label{sec:edge_cases}

\subsection{Failure Modes at Low Cardinality}
\label{sec:low_card_failure}

The \texttt{low\_card} config (18/12 rows/entity, 100 entities, 10\% overlap) exposes fundamental failure modes:

\subsubsection{Independence Assumption Violation}
The \fellegisunter{} model assumes conditional independence of comparisons given match status.
At low cardinality, each \globalentity{} has few records ($m,n \approx 5$).
The local pairwise comparisons between records of the \textit{same} \globalentity{} are highly correlated—they share the same latent true values corrupted independently.
Formally, for $r_{A,i}, r_{B,j} \sim e$:
\begin{equation}
	\text{Cov}(\gamma(r_{A,i}, r_{B,j}), \gamma(r_{A,i'}, r_{B,j'})) > 0 \quad \text{for } i \neq i' \text{ or } j \neq j'
\end{equation}
This inflates the effective evidence: the product of likelihood ratios overcounts the information.
Result: match weights are overconfident (too high for matches, too low for non-matches), but the \textit{ranking} may still be reasonable.

\subsubsection{Combinatorial Explosion in Comparison Space}
For a global pair with $m$ and $n$ records, there are $m \times n$ local pairs.
When $m$,$n$ are large (e.g., $18 \times 12 = 216$ pairs per global entity), the Cartesian product explodes.
Blocking must be aggressive, but aggressive blocking loses recall.

\subsubsection{Sparse Parameter Estimation}
\emalgo{} estimates $m_k(\gamma)$ from pairs in the blocking rule.
At low cardinality, the number of true matches in any blocking rule is small ($O(N \rho)$ with small $N$).
Parameter estimates have high variance; \emalgo{} may converge to unstable values or fail to converge.

\subsubsection{Aggregation robustness}
As we see in \ref{sec:exp_aggregation}, single scalar aggregation functions handle most cardinality regimes robustly. At regimes of extremely low cardinality, however, this approach breaks - record pairs are degenerate, and function scores are as a result, noisy.

% \subsection{Proposed Mitigation: Profile-Based Probabilistic Matching}
% \label{sec:profile_splink_proposal}

% \textbf{Core idea}: Instead of aggregating scalar scores, collect \textit{orthogonal similarity scores} as features and apply \fellegisunter{} (\splink{}) on the profile vectors.

% \subsubsection{Step 1: Collect Orthogonal Similarities}
% For each global entity pair $(e_A, e_B)$, compute a feature vector $\mathbf{x} \in \mathbb{R}^D$:
% \begin{align}
% 	x_1 & = \max_{i,j} p_{ij}                                                   & \text{(max local prob)}  \\
% 	x_2 & = \frac{1}{mn} \sum_{i,j} p_{ij}                                      & \text{(mean local prob)} \\
% 	x_3 & = 1 - \prod_{i,j} (1-p_{ij})                                          & \text{(noisy-OR)}        \\
% 	x_4 & = \text{soft Jaccard}(\text{title\_tokens}_A, \text{title\_tokens}_B)                            \\
% 	x_5 & = \text{Simpson}(\text{track\_nos}_A, \text{track\_nos}_B)                                       \\
% 	x_6 & = \text{range overlap}(\text{length}_A, \text{length}_B)                                         \\
% 	x_7 & = \text{year overlap count}                                                                      \\
% 	x_8 & = \text{mag-cosine}(\text{language\_vec}_A, \text{language\_vec}_B)                              \\
% 		    & \vdots
% \end{align}
% These $D$ features are designed to be \textit{as orthogonal as possible}—each captures a different matching signal.

% \subsubsection{Step 2: Adjust for Correlation}
% The features $x_d$ are not conditionally independent given match status.
% We estimate their correlation structure on a held-out set (or via bootstrap on the unsupervised predictions):
% \begin{equation}
% 	\mathbf{\Sigma}_M = \text{Cov}(\mathbf{x} \mid M), \quad \mathbf{\Sigma}_U = \text{Cov}(\mathbf{x} \mid U)
% \end{equation}
% Then the profile-based \fellegisunter{} uses a multivariate Gaussian (or copula) likelihood instead of the product of univariate likelihoods:
% \begin{equation}
% 	\frac{p(\mathbf{x} \mid M)}{p(\mathbf{x} \mid U)} = \frac{\mathcal{N}(\mathbf{x}; \boldsymbol{\mu}_M, \mathbf{\Sigma}_M)}{\mathcal{N}(\mathbf{x}; \boldsymbol{\mu}_U, \mathbf{\Sigma}_U)}
% \end{equation}
% Or, pragmatically: use \splink{}'s \texttt{CustomComparison} with binned levels on each $x_d$, and accept the residual dependence (as in the primary layer).

% \subsubsection{Step 3: Profile-Based \splink{} Model}
% Apply a \splink{} model on the global-pair feature vectors:
% \begin{itemize}
% 	\item Input: DataFrame with one row per global entity pair, columns = $x_1 \dots x_D$.
% 	\item Blocking: Use $x_1$ (max prob) or $x_3$ (noisy-OR) as blocking key.
% 	\item \emalgo{} estimation: Same procedure; now each ``record'' is a global entity pair.
% 	\item Output: Calibrated global match probability.
% \end{itemize}

% \subsubsection{Preliminary Results on \texttt{low\_card}}
% \label{sec:low_card_prelim}
% Experiments in \texttt{splink\_low\_card.ipynb} and \texttt{splink\_low\_card\_copy.ipynb}:
% \begin{itemize}
% 	\item Primary layer on \texttt{low\_card}: Match weight separation degraded; $F_1$ drops to \textit{TBD}.
% 	\item Profile construction: Aggregator runs on 18/12-record groups; vectors are sparse but informative.
% 	\item Profile-based \splink{}: \textit{TBD: Insert results from notebooks}.
% 	\item Comparison: Profile-based vs best scalar aggregation (max).
% \end{itemize}

% \begin{table}[htbp]
% 	\centering
% 	\caption{Cardinality Extremes: Method Comparison}
% 	\label{tab:cardinality_comparison}
% 	\begin{tabular}{lcccc}
% 		\toprule
% 		\textbf{Method}        & \textbf{Low Card (5/5)} & \textbf{Standard (100/120)} & \textbf{High Card (200/200)} \\
% 		\midrule
% 		Primary layer (threshold) & \textit{TBD}            & \textit{TBD}                & \textit{TBD}                 \\
% 		Naive Sum              & \textit{TBD}            & \textit{TBD}                & \textit{TBD}                 \\
% 		Normalised Mean        & \textit{TBD}            & \textit{TBD}                & \textit{TBD}                 \\
% 		Noisy-OR               & \textit{TBD}            & \textit{TBD}                & \textit{TBD}                 \\
% 		Max                    & \textit{TBD}            & \textit{TBD}                & \textit{TBD}                 \\
% 		Vector-Space (primary layer) & \textit{TBD}            & \textit{TBD}                & \textit{TBD}                 \\
% 		Profile-based \splink{}      & \textit{TBD}            & \textit{TBD}                & \textit{TBD}                 \\
% 		\bottomrule
% 	\end{tabular}
% \end{table}

% \subsection{When Splink + Aggregation Fails: Decision Boundary}
% \label{sec:failure_boundary}

% The hybrid approach breaks down when:
% \begin{enumerate}
% 	\item \textbf{Cardinality $> 50$ per entity}: Blocking cannot reduce pairs sufficiently; \emalgo{} estimation becomes unstable.
% 	\item \textbf{Overlap $< 1\%$}: $\pi$ too small; \emalgo{} converges to trivial solution (everything non-match).
% 	\item \textbf{Attribute overlap $< 2$ fields}: Insufficient signal for separation.
% 	\item \textbf{Within-dataset duplicates not resolved}: If \texttt{local\_source\_id} is not stable, the local-to-global decomposition is invalid.
% \end{enumerate}
% In these regimes, alternative approaches (graph-based collective resolution, deep blocking, active learning) are needed.
% \end{document}
